\chapter{Introduction}\label{chap:intro}

\section{Background and Motivation}\label{sec:intro_background}

The way people organise their work has changed fundamentally over the last decade.
The proliferation of distributed teams, freelance arrangements, and hybrid work models
has driven enormous growth in the market for digital task-management and project-collaboration
software~\cite{allen2001gtd}.
Tools such as Trello, Asana, Jira, and Notion have become everyday staples in software
development teams, enabling members to track progress, share context, and coordinate
across time zones without the overhead of physical meetings.

In practice, many organisations---universities, government institutions, software companies, and large enterprises---still communicate task assignments and announcements through social-media platforms or messaging applications. As message volume grows, important tasks become buried among thousands of notifications, making them difficult to track and increasing the likelihood that critical deadlines or events are overlooked. Managers lack a structured mechanism to monitor progress, assess performance, and generate objective reports on task completion and delays. Without a dedicated system, generating such reports is time-consuming and often impractical.

TaskFlow addresses these challenges by centralising task assignment, progress monitoring, reporting, collaboration, and intelligent notifications within a single desktop environment, ensuring that responsibilities are visible, traceable, and never lost in an overflowing chat feed.

Despite this progress, the dominant generation of collaboration tools shares a critical
architectural assumption: \emph{a reliable, always-available internet connection}.
Every create, read, update, and delete operation is routed through a remote cloud server.
When connectivity degrades---as it does on mobile hotspots, in aeroplanes, in rural areas,
or during network maintenance---these applications either display error banners or silently
discard user input.
For a professional whose productivity depends on continuous access to their task list,
such fragility is unacceptable.

A second limitation is the growing gap between the capabilities of modern Large Language
Models (LLMs) and the AI integrations available in mainstream productivity tools.
Although several task managers have begun experimenting with AI-generated summaries,
none provides a fully embedded, context-aware AI assistant that understands the user's
current project state and can assist with both natural-language planning and technical
code generation within the same interface~\cite{jiang2023mistral}.

These observations---the limitations of ad-hoc communication channels for task management, the unreliability of purely cloud-centric architectures, and the
absence of deep AI integration---form the motivation for \textit{TaskFlow}: a
\emph{Smart Task \& Quality Management System} designed from the ground up as an
offline-first desktop application with an embedded multi-mode AI assistant.


\section{Problem Statement}\label{sec:intro_problem}

The central problem addressed by this project can be stated as follows:

\begin{quote}
\textit{Existing task-management tools treat internet connectivity as a prerequisite
rather than an enhancement.
They provide no graceful degradation path, no local data store, and no mechanism to
synchronise divergent offline edits once connectivity is restored.
Furthermore, their AI capabilities, where present, are superficial and disconnected
from the user's live project data.}
\end{quote}

More concretely, a review of widely-used tools reveals the following gaps.

\begin{enumerate}
    \item \textbf{No offline persistence.}
          Cloud-first tools store all state remotely. A user working without connectivity
          cannot create tasks, update statuses, or leave comments; changes attempted offline
          are either silently lost or surfaced as errors.

    \item \textbf{No transparent synchronisation.}
          Even tools that offer limited offline caching typically give users no visibility into
          which changes are pending, in progress, or confirmed. The sync lifecycle is a
          black box, eroding trust in the accuracy of the displayed data.

    \item \textbf{Fragmented toolchains.}
          Developers and project managers routinely switch between a task manager, a code
          assistant, a calendar, a messaging client, and a notification centre.
          There is no single, cohesive desktop application that integrates all of these
          capabilities into one consistent experience.

    \item \textbf{Shallow AI integration.}
          AI features in existing tools are limited to simple text generation (e.g.,
          auto-writing a task description). None provides a conversational assistant
          capable of code generation, document analysis, or multimodal input---all
          within the context of the user's own project data.
\end{enumerate}

\section{Aims and Objectives}\label{sec:intro_aims}

The primary aim of this project is to design, implement, and evaluate a cross-platform
desktop task-management application that operates reliably regardless of network
availability, while embedding a powerful AI assistant and a full suite of collaboration
features.

The following specific objectives guide the work:

\begin{enumerate}
    \item \textbf{Offline-first operation.}
          Every core feature---task creation and editing, project management, calendar
          events, notifications, and messaging---must function fully without an active
          internet connection.
Data must be persisted locally in a SQLite database for performance, with MongoDB
serving as the single source of truth for synchronisable data across devices.

    \item \textbf{Reliable, transparent synchronisation.}
          When connectivity is available, locally generated changes must be propagated to
          a MongoDB relay database using a transactional outbox pattern that guarantees
          at-least-once delivery.
          The user must be able to see the current connectivity status, the count of
          pending operations, and the real-time progress of any active sync pass through
          a dedicated UI indicator.

    \item \textbf{Integrated multi-mode AI assistant.}
          A conversational AI assistant backed by the Mistral API must be embedded
          directly in the application.
          It must support three operational modes: a general-purpose chat mode, a
          code-generation mode powered by Codestral, and a document-analysis mode
          capable of processing uploaded images and PDF files via optical character
          recognition (OCR).
          Responses must be streamed to the UI via Server-Sent Events (SSE) to provide
          a responsive, real-time interaction feel.

    \item \textbf{Five complementary task visualisation views.}
          The task module must offer five distinct views catering to different cognitive
          styles and workflows: a smart-grouped default list, a Kanban board with
          drag-and-drop, a tabular data-grid, a Gantt-style timeline, and a calendar
          time-grid.

    \item \textbf{Cross-platform desktop distribution.}
          The completed application must be packaged as a native installer for
          Windows (NSIS and MSI) using Tauri~v2, bundling both the .NET backend and
          the React frontend into a single self-contained executable that requires
          no external runtime installation by the end user.
          The Tauri framework was chosen for its significantly smaller installer size
          and faster cold-start performance, achieved by leveraging the operating
          system's native WebView rather than bundling a browser engine.

    \item \textbf{Comprehensive collaboration feature set.}
          The system must include team management, one-to-one and group messaging with
          file attachments, real-time notifications via SignalR, calendar events,
          a dashboard with productivity statistics, and user-profile settings including
          security and preference controls.
\end{enumerate}

\section{Scope and Limitations}\label{sec:intro_scope}

\subsection*{In Scope}

The following capabilities were designed and implemented as part of this project:

\begin{itemize}
    \item Full user authentication with local-identity middleware, bcrypt password hashing, and
          local file-based password-reset workflow.
    \item Task management with priority levels (Low, Medium, High), status workflow
          (To Do, In Progress, Review, Completed, Overdue), due-date tracking, starring,
          and assignment.
    \item Five task visualisation views: Default, Kanban, Table, Gantt, and Calendar.
    \item Project and team management with role-based membership and team invitations
          synchronised through MongoDB.
    \item One-to-one direct messaging and group chats, both supporting file attachments
          and soft-delete.
    \item In-application and email reminders triggered by configurable fire-at times.
    \item Automated due-date warnings fired 24~hours and 1~hour before deadline.
    \item Real-time push notifications delivered over SignalR.
    \item An AI chatbot with three modes, SSE streaming, and conversation history.
    \item Calendar events with a weekly time-grid view.
    \item A dashboard with live statistics, activity feeds, and a calendar widget.
    \item Offline-first SQLite persistence and an outbox-based MongoDB synchronisation
          subsystem.
    \item Cross-platform desktop packaging using Tauri~v2, producing a
          self-contained Windows installer with a significantly reduced footprint.
\end{itemize}

\subsection*{Limitations}

The following items were explicitly excluded from the current implementation:

\begin{itemize}
    \item \textbf{Mobile clients.}
          iOS and Android applications were not developed; the system targets desktop
          platforms only.
    \item \textbf{Multi-user conflict resolution.}
          The synchronisation subsystem uses a last-write-wins strategy.
          Collaborative real-time editing with Conflict-free Replicated Data Types
          (CRDTs)~\cite{shapiro2011crdt} is a known future enhancement.
    \item \textbf{Hierarchical sub-tasks.}
          A sub-task UI scaffold is present in the Table view but the backend persistence
          layer for sub-tasks was not implemented in this iteration.
    \item \textbf{SaaS deployment and billing.}
          The application is designed for single-user or small-team use on a local machine;
          multi-tenant cloud hosting and subscription management are out of scope.
    \item \textbf{Externally managed AI credentials.}
          The Mistral API key is currently supplied at build time.
          Secure external key management (e.g., system keychain or environment-variable
          injection) is identified as a known improvement area in
          Chapter~\ref{chap:concl}.
\end{itemize}

\section{Thesis Overview}\label{sec:intro_overview}

The remainder of this thesis is structured as follows.

\textbf{Chapter~\ref{chap:background} --- Literature Review and Background} surveys
the academic and technical context of the work, covering task-management paradigms,
offline-first architecture, hybrid database design, real-time communication protocols,
authentication, and the state of LLM-based AI assistants.

\textbf{Chapter~\ref{chap:Methodology} --- Methodology and System Design} presents the
full system architecture, the rationale behind key design decisions, the data model, and
a detailed walkthrough of the backend, frontend, and desktop-packaging implementations.

\textbf{Chapter~\ref{chap:results} --- Results and System Evaluation} demonstrates the
implemented features through functional descriptions and user-interface walkthroughs,
assessing the system against each objective stated in Section~\ref{sec:intro_aims}.

\textbf{Chapter~\ref{chap:concl} --- Conclusion} reflects on the outcomes of the
project, summarises its technical contributions, and acknowledges known limitations.

\textbf{Chapter~\ref{chap:todo} --- Future Work} outlines a prioritised roadmap of
short-term improvements and long-term enhancements that could extend the system.

\textbf{Appendices} provide supplementary technical reference material including the
REST API endpoint catalogue, the complete database schema, and the application
configuration reference.
